\chapter{Need, task e storyboard}
\section{Need}
In seguito all’analisi delle risposte alle interviste e delle risposte ai questionari, i principali need che sono emersi sono i seguenti (in ordine di decrescenza e importanza):
\begin{enumerate}
    \item necessità di conoscere lo stato attuale di un’infrastruttura (affollamento, pulizia, sicurezza ecc.)
    \item necessità di conoscere eventi e iniziative e le loro informazioni
    \item necessità di trovare e interagire con persone con interessi simili
    \item necessità di conoscere infrastrutture verdi nelle vicinanze e le loro informazioni
    \item necessità di conoscere i servizi offerti e le loro informazioni di un’infrastruttura verde
\end{enumerate}
\section{Task}
Per quanto riguarda i task, abbiamo trovato i seguenti task e sotto-task:
\begin{enumerate}
    \item[1.] L’utente può conoscere lo stato attuale di un’infrastruttura da lui scelta. 

    L’utente:
    \begin{enumerate}
        \item Seleziona un’infrastruttura verde
        \item Visualizza le informazioni relative allo stato attuale dell’infrastruttura
    \end{enumerate}

    \item[1.1.] L’utente aggiorna lo stato di un’infrastruttura verde (inserimento/modifica).
    
	L’utente:
    \begin{enumerate}
        \item Cerca un’infrastruttura verde
        \item Seleziona un’infrastruttura verde
        \item Aggiorna i parametri attuali dell’infrastruttura selezionata
        \item Inserisce una valutazione sui vari aspetti dell’infrastruttura verde (affollamento, pulizia, ...)
    \end{enumerate}

    \item[2.] L’utente può conoscere gli eventi e iniziative in una infrastruttura verde:
    
    L’utente:
    \begin{enumerate}
        \item Seleziona un’infrastruttura verde da lui cercata
        \item Inserisce eventuali filtri sugli eventi (date, orari, ...)
        \item Visualizza le informazioni degli eventi
    \end{enumerate}

    \item[2.1.] L’utente può visualizzare gli eventi e iniziative nelle infrastrutture verdi con eventuali filtri:
    
    L’utente:
    \begin{enumerate}
        \item Inserisce eventuali filtri sugli eventi (date, orari, ...)
        \item Inserisce eventuali filtri (distanza a piedi, tipo di infrastruttura, ...)
        \item Visualizza le informazioni sulle infrastrutture verdi e sugli eventi
    \end{enumerate}
    
    \item[3.] L’utente può interagire con altre persone nelle infrastrutture verdi
    
    L’utente:
    \begin{enumerate}
        \item Imposta i filtri di ricerca (attività, distanza, ...)
        \item Visualizza tutti gli altri utenti che hanno condiviso la propria posizione e attività
        \item L’utente contatta la persona selezionata per partecipare all’attività
    \end{enumerate}

    \item[3.1.] L’utente può condividere la propria attività in una infrastruttura verde (inserimento/modifica)
    
    L’utente:
    \begin{enumerate}
        \item inserisce la posizione attuale
        \item inserisce l’attività che sta svolgendo
    \end{enumerate}
    
    \item[4.] L’utente può cercare infrastrutture verdi nelle vicinanze.

    L’utente:
    \begin{enumerate}
        \item seleziona la propria posizione
        \item inserisce eventuali filtri (distanza a piedi, tipo di infrastruttura, ...)
        \item seleziona l’infrastruttura verde desiderata
    \end{enumerate}
    
    \item[5.] L’utente può conoscere i servizi offerti di un’infrastruttura verde

    L’utente:
    \begin{enumerate}
        \item Seleziona un’infrastruttura verde
        \item Visualizza le informazioni dei servizi disponibili
    \end{enumerate}
\end{enumerate}

\section{Storyboard}
Sono stati realizzati in totale 5 storyboards rappresentanti gli svolgimenti dei principali 5 task.
Per esaminare i disegni relativi agli storyboards, utilizzare \href{https://drive.google.com/drive/u/0/folders/13NUPKKbf_ISkTarvQYw04d_vPzp9hNU0}{il seguente link} alla cartella drive del progetto.